% ==========================================
% PREÁMBULO Y CONFIGURACIÓN DE PAQUETES
% ==========================================

\usepackage[spanish,es-tabla]{babel}
\usepackage[utf8]{inputenc}
\usepackage[margin=2.5cm]{geometry}

% Paquetes matemáticos y de física
\usepackage{amsmath,amssymb,amsfonts,amsthm}
\usepackage{physics}
\usepackage{siunitx} % Para magnitudes físicas y unidades SI

% -----------------------------------------------------------------
% DESACTIVACIÓN DEFINITIVA DEL WARNING SIUNITX + PHYSICS
% -----------------------------------------------------------------
% 1. Desactiva el mensaje de advertencia interno de LaTeX3
\ExplSyntaxOn
\msg_redirect_name:nnn { siunitx } { physics-pkg } { none }
\ExplSyntaxOff

% 2. Asigna el comando \qty a siunitx (para escribir ej: \qty{5}{\volt})
\AtBeginDocument{\RenewCommandCopy\qty\SI}
% -----------------------------------------------------------------

% Formato de listas e imágenes
\usepackage{graphicx}
\usepackage{booktabs}
\usepackage{xcolor}
\usepackage{hyperref}
\usepackage{enumitem}
\usepackage{capt-of}

% Cajas de texto personalizadas
\usepackage[most]{tcolorbox}
\tcbuselibrary{breakable}

\newcounter{problemacontador}[chapter]

% Definición de colores principales
\definecolor{maincolor}{RGB}{0, 51, 102}      % Azul oscuro académico
\definecolor{accentcolor}{RGB}{0, 128, 128}   % Verde azulado / Teal
\definecolor{boxbg}{RGB}{245, 247, 250}       % Fondo claro para cajas

% Configuración de hyperref
\hypersetup{
    colorlinks=true,
    linkcolor=maincolor,
    filecolor=magenta,      
    urlcolor=cyan,
    citecolor=accentcolor
}

% Entornos para conceptos clave y fórmulas destacadas
\newtcolorbox{concepto}[1][]{
  colback=boxbg,
  colframe=maincolor,
  fonttitle=\bfseries,
  title={#1},
  arc=2mm,
  boxrule=1pt
}

\newtcolorbox{formula}[1][]{
  colback=boxbg,
  colframe=accentcolor,
  fonttitle=\bfseries,
  title={Fórmula clave: #1},
  arc=2mm,
  boxrule=1pt
}

\newtcolorbox{problema}[1][]{
  enhanced,
  breakable,
  colback=boxbg,
  colframe=maincolor,
  fonttitle=\bfseries,
  coltitle=white,
  arc=2mm,
  boxrule=1pt,
  before upper={\refstepcounter{problemacontador}},
  title={Problema~\theproblemacontador\ifstrempty{#1}{}{: #1}},
  lower separated=true
}

% Configuración de unidades físicas
\sisetup{
  per-mode = symbol,
  separate-uncertainty = true
}